% =============================================================================
% Parameter Calibration Section
% Meta-Prediction Theory Paper
% =============================================================================

\section{Parameter Calibration and Model Validation}

This section describes the calibration of the computational model parameters using empirical data and validates the model predictions against observed patterns.

\subsection{Model Overview}

The Meta-Prediction Model describes the process by which respondents generate desired answers through recursive meta-cognitive processes. The model includes five key parameters:

\begin{enumerate}
    \item \textbf{Memory Accuracy ($\alpha$)}: Accuracy of recalling actual experiences
    \item \textbf{Self-Awareness ($\sigma$)}: Clarity of self-perception
    \item \textbf{Ideal Self-Strength ($\delta$)}: Strength of ideal self-concept
    \item \textbf{Norm Awareness ($\nu$)}: Awareness of social expectations
    \item \textbf{Meta-Cognitive Capacity ($\mu$)}: Ability to monitor and regulate responses
\end{enumerate}

\subsection{Parameter Estimation from Empirical Data}

\subsubsection{Data Source}

We used the Temporal Dynamics dataset ($N = 24,292$) to estimate model parameters. The dataset includes:
\begin{itemize}
    \item PHQ-9 responses (9 items)
    \item GAD-7 responses (7 items)
    \item Response times for each item
\end{itemize}

\subsubsection{Estimation Method}

Parameters were estimated using ordinary least squares (OLS) regression with the Temporal Dynamics dataset ($N = 24{,}292$). The five model parameters were derived from the structural path coefficients using a logistic transformation to bound them in the $[0, 1]$ range:

\subsubsection{Estimated Parameters}

Table \ref{tab:parameters} presents the estimated parameter values.

\begin{table}[h]
\centering
\begin{tabular}{lccc}
\toprule
\textbf{Parameter} & \textbf{Symbol} & \textbf{Estimate} & \textbf{95\% CI} \\
\midrule
Memory Accuracy & $\alpha$ & 0.61 & [0.61, 0.62] \\
Self-Awareness & $\sigma$ & 0.74 & [0.73, 0.75] \\
Ideal Self-Strength & $\delta$ & 1.00 & [1.00, 1.00] \\
Norm Awareness & $\nu$ & 0.50 & [0.50, 0.50] \\
Meta-Cognitive Capacity & $\mu$ & 0.53 & [0.51, 0.54] \\
\bottomrule
\end{tabular}
\caption{Estimated Model Parameters}
\label{tab:parameters}
\end{table}

\subsection{Model Validation}

\subsubsection{Validation Approach}

We validated the model by comparing its predictions against observed patterns in the empirical data:

\begin{enumerate}
    \item \textbf{Response time patterns}: Model-predicted response times vs. observed
    \item \textbf{Score distributions}: Model-predicted score distributions vs. observed
    \item \textbf{Correlation patterns}: Model-predicted correlations vs. observed
\end{enumerate}

\subsubsection{Response Time Predictions}

The model predicts that response times will increase with:
\begin{itemize}
    \item Item difficulty (higher scores)
    \item Meta-cognitive monitoring (greater deliberation)
    \item Scale awareness (greater impression management)
\end{itemize}

Figure \ref{fig:rt_predictions} shows the comparison between model-predicted and observed response times.

\begin{figure}[h]
\centering
\includegraphics[width=0.8\textwidth]{rt_predictions.png}
\caption{Model-Predicted vs. Observed Response Times}
\label{fig:rt_predictions}
\end{figure}

\subsubsection{Score Distribution Predictions}

The model predicts that score distributions will be:
\begin{itemize}
    \item Right-skewed (most respondents score low)
    \item More variable for respondents with higher meta-cognitive capacity
    \item Less variable for respondents with higher norm awareness
\end{itemize}

Table \ref{tab:distribution} compares model-predicted and observed distribution statistics.

\begin{table}[h]
\centering
\begin{tabular}{lcc}
\toprule
\textbf{Statistic} & \textbf{Observed} & \textbf{Predicted} \\
\midrule
Mean & 3.27 & 3.15 \\
SD & 3.71 & 3.58 \\
Skewness & 1.64 & 1.38 \\
Kurtosis & 4.05 & 2.08 \\
\bottomrule
\end{tabular}
\caption{Comparison of Observed and Predicted Distributions}
\label{tab:distribution}
\end{table}

\subsubsection{Correlation Pattern Predictions}

The model predicts the following correlation patterns:

\begin{itemize}
    \item PHQ-9 $\leftrightarrow$ GAD-7: $r = 0.76$ (observed: 0.77)
    \item PHQ-9 $\leftrightarrow$ Response Time: $r = 0.09$ (observed: 0.09)
    \item Scale Awareness $\leftrightarrow$ PHQ-9: $r = -0.50$ (observed: -0.50)
\end{itemize}

Table \ref{tab:correlations_model} presents the full comparison.

\begin{table}[h]
\centering
\begin{tabular}{lcc}
\toprule
\textbf{Correlation} & \textbf{Observed} & \textbf{Predicted} \\
\midrule
PHQ-9 $\leftrightarrow$ GAD-7 & 0.766 & 0.758 \\
PHQ-9 $\leftrightarrow$ RT Mean & 0.091 & 0.088 \\
PHQ-9 $\leftrightarrow$ SA & -0.498 & -0.492 \\
PHQ-9 $\leftrightarrow$ MM & -0.132 & -0.128 \\
PHQ-9 $\leftrightarrow$ DEL & -0.296 & -0.289 \\
\bottomrule
\end{tabular}
\caption{Comparison of Observed and Predicted Correlations}
\label{tab:correlations_model}
\end{table}

\subsection{Model Comparison}

We compared the Meta-Prediction Model against two alternative models:

\begin{enumerate}
    \item \textbf{Null Model}: No meta-cognitive processes (simple responding)
    \item \textbf{Concurrent-Validity Model}: Only GAD-7 as predictor (baseline concurrent validity)
\end{enumerate}

Table \ref{tab:model_comparison} presents the model comparison results.

\begin{table}[h]
\centering
\begin{tabular}{lcc}
\toprule
\textbf{Model} & \textbf{AIC} & \textbf{BIC} \\
\midrule
Null Model & 132,614 & 132,622 \\
Concurrent-Validity Model & 111,131 & 111,147 \\
Meta-Prediction Model & 92,143 & 92,191 \\
\bottomrule
\end{tabular}
\caption{Model Comparison Results}
\label{tab:model_comparison}
\end{table}

The Meta-Prediction Model showed the best fit to the data (lowest AIC and BIC), supporting the theoretical framework.

\subsection{Sensitivity Analysis}

We conducted sensitivity analyses to assess the robustness of parameter estimates:

\begin{enumerate}
    \item \textbf{Bootstrap analysis}: 1,000 bootstrap samples
    \item \textbf{Cross-validation}: 5-fold cross-validation
\end{enumerate}

\subsubsection{Bootstrap Results}

Table \ref{tab:bootstrap} presents the bootstrap results for parameter estimates.

\begin{table}[h]
\centering
\begin{tabular}{lccc}
\toprule
\textbf{Parameter} & \textbf{Estimate} & \textbf{Bootstrap SD} & \textbf{95\% CI} \\
\midrule
$\alpha$ & 0.61 & 0.004 & [0.60, 0.62] \\
$\sigma$ & 0.74 & 0.016 & [0.71, 0.77] \\
$\delta$ & 1.00 & 0.000 & [1.00, 1.00] \\
$\nu$ & 0.50 & 0.001 & [0.50, 0.50] \\
$\mu$ & 0.53 & 0.007 & [0.52, 0.55] \\
\bottomrule
\end{tabular}
\caption{Bootstrap Results for Parameter Estimates}
\label{tab:bootstrap}
\end{table}

\subsubsection{Cross-Validation Results}

The 5-fold cross-validation showed:
\begin{itemize}
    \item Mean AIC: 73,714 (SD = 138)
    \item Mean BIC: 73,761 (SD = 138)
    \item Stable parameter estimates across folds
\end{itemize}

\subsection{Implications for Theory}

The parameter calibration results have several implications for the Meta-Prediction Theory:

\begin{enumerate}
    \item \textbf{Self-Awareness is highest} ($\sigma = 0.74$): The cross-scale consistency variable shows the strongest association with PHQ-9 scores, though this partly reflects the circular relationship between this variable and the outcome measure (see Limitations).
    
    \item \textbf{Deliberation is strongest} ($\delta = 1.00$): The inverse-extreme-responding variable shows the largest effect, suggesting that response style---specifically, the tendency to avoid extreme responses---is the dominant predictor among the meta-cognitive proxies.
    
    \item \textbf{Meta-Cognitive Capacity is moderate} ($\mu = 0.53$): RT variability shows a modest negative association with PHQ-9 scores, consistent with the interpretation that greater response uncertainty is associated with lower reported symptoms.
    
    \item \textbf{Memory Accuracy is moderate} ($\alpha = 0.61$): GAD-7 shows a strong positive association with PHQ-9, consistent with substantial shared variance between depression and anxiety measures.
    
    \item \textbf{Norm Awareness is negligible} ($\nu = 0.50$): Response time mean shows a very small positive coefficient, suggesting that average response speed has minimal predictive value beyond the other variables.
\end{enumerate}

\subsection{Conclusion}

The parameter calibration and model validation provide strong support for the Meta-Prediction Theory. The model fits the empirical data better than alternative models, and the parameter estimates are consistent with theoretical expectations and robust across validation approaches.

\end{document}
